\begin{abstract}
The owner of an existing blockchain address wants to make one move, go offline and have a committee build and sign one spend when a condition on the chain holds. Published designs hand the committee the key, so that a threshold of nodes can sign anything from then on, or they keep the owner in the protocol. We define the One-Move Spend problem and an ideal functionality $\FOMS$ for it, and we give a protocol for it, \emph{line dealing}, in two modes. A Schnorr or ECDSA signature is one point on a line whose two coefficients the owner knows. The owner's single message is a publicly checkable Shamir dealing, on the signature curve, of \emph{sealed points} or of \emph{open lines}. A sealed point is a finished signature on one transaction of a fee grid whose members all spend one pinned input. An open line is both coefficients, under commitments whose constant terms are forced to the generator and the address key. The committee runs no distributed key generation, no presignature protocol and no secret multiplication. Sealed mode holds a fully corrupt committee to the owner's conflicting transactions, of which the chain executes at most one, and it never exposes the key. Open mode lets the committee build the transaction from chain state, lock its digest with a quorum certificate before any share is used and sign with one linear partial signature per node. Its one-spend limit holds while at most $f$ of $n \ge 3f+1$ nodes are Byzantine. We prove that open mode cannot have the sealed form of confinement: when an honest committee would sign either of two messages that later chain state selects, no one-move protocol confines a qualified coalition to one signature. Two signatures need not be two executed spends, since both may spend the pin; open mode's one-spend limit rests on an honest threshold because $t$ nodes reconstruct the key it deals. Open ECDSA is reduced to the theorem of Groth and Shoup on ECDSA with presignatures. Open Schnorr is proved under the algebraic one-more discrete logarithm assumption, and the share encryption under interactive CDH. Resharing to a new committee, adaptive corruption and a proof that the protocol realizes $\FOMS$ are open. The algebra and the counting arguments are machine-checked in Lean~4 (103 theorems), and the board rules are explored by randomized simulation of a Quint model. At $n = 9$ a Rust prototype spends about \nineNodeMicrosLow{}\us{} of CPU per node on a request, about \nineReqPerCoreLow{} requests per second per core.
\end{abstract}
